\section{Conclusion and Future Work}
\label{sec:conclusion}

In this paper, we have presented a novel second-order geometric framework for adaptive model merging, addressing a fundamental vulnerability in unsupervised test-time coefficient optimization. We formalize this vulnerability as the \textbf{Overfitting-Optimizer Paradox}: in the absence of constraints, test-time adaptation of layer-wise merging coefficients fits local transductive stream noise, triggering high-frequency spatial coefficient oscillations and catastrophic internal representation collapse. 

To resolve this paradox, we introduced \textbf{Riemannian Curvature-Regularized Test-Time Model Merging (RCR-Merge)}. By modeling the network parameter space as a Riemannian manifold where distance is locally scaled by the trace of the base model's diagonal Fisher Information Matrix (FIM), RCR-Merge regularizes spatial total variation dynamically across network depth. We proved theoretically (Lemma~\ref{lem:coordinate_barrier} and Theorem~\ref{thm:representation_drift}) that scaling adjacent-layer penalties with the geometric mean of their curvatures construct an analytical barrier, shielding highly sensitive bottleneck layers from harmful coefficient fluctuations while preserving optimization adaptability in robust representation layers. 

Our empirical evaluations across 30 seeds on the Coupled Model II Landscape validate our theoretical claims. RCR-Merge successfully resolves the overfitting collapse, outperforming the static Uniform Baseline by \textbf{+3.06\% absolute}, unconstrained AdaMerging by \textbf{+10.33\% absolute}, and flat Total Variation regularization by \textbf{+5.01\% absolute}, while demonstrating unparalleled trajectory stability and variance reduction compared to rigid subspace methods. 

\paragraph{Future Work.}
There are several promising avenues for future research. First, we plan to implement and evaluate our formulated Kronecker-factored Fisher Information Matrix (K-FAC) extension on actual large-scale transformer architectures. By capturing intra-layer anisotropic parameter correlations through lightweight, low-dimensional input and pre-activation gradient covariance matrices (such as $A^{(l)}$ and $G^{(l)}$), K-FAC RCR-Merge can protect attention channels and MLP sub-blocks individually with near-zero computational overhead. Second, we plan to extend RCR-Merge to modern 7B and 70B parameter large language models (LLMs) and autoregressive multimodal architectures under long-term non-stationary streams. Third, we aim to derive formal generalization bounds for adaptive model merging using localized Rademacher complexity and Wasserstein distance on Riemannian parameter manifolds. Finally, we believe that integrating dynamic, online updates to the base curvature estimates during test-time adaptation could further improve robust adaptability under severe out-of-distribution shifts. Overall, our work demonstrates that grounding optimization trajectories in the second-order geometric landscape of pre-trained models is a powerful paradigm for robust, adaptive, and scalable multi-task model merging.